\documentclass[a4paper,11pt,fleqn]{article}\usepackage[]{graphicx}\usepackage[]{color}
%% maxwidth is the original width if it is less than linewidth
%% otherwise use linewidth (to make sure the graphics do not exceed the margin)
\makeatletter
\def\maxwidth{ %
  \ifdim\Gin@nat@width>\linewidth
    \linewidth
  \else
    \Gin@nat@width
  \fi
}
\makeatother

\definecolor{fgcolor}{rgb}{0, 0, 0}
\newcommand{\hlnum}[1]{\textcolor[rgb]{0,0,0}{#1}}%
\newcommand{\hlstr}[1]{\textcolor[rgb]{0,0,0}{#1}}%
\newcommand{\hlcom}[1]{\textcolor[rgb]{0.4,0.4,0.4}{\textit{#1}}}%
\newcommand{\hlopt}[1]{\textcolor[rgb]{0,0,0}{\textbf{#1}}}%
\newcommand{\hlstd}[1]{\textcolor[rgb]{0,0,0}{#1}}%
\newcommand{\hlkwa}[1]{\textcolor[rgb]{0,0,0}{\textbf{#1}}}%
\newcommand{\hlkwb}[1]{\textcolor[rgb]{0,0,0}{\textbf{#1}}}%
\newcommand{\hlkwc}[1]{\textcolor[rgb]{0,0,0}{\textbf{#1}}}%
\newcommand{\hlkwd}[1]{\textcolor[rgb]{0,0,0}{\textbf{#1}}}%
\let\hlipl\hlkwb

\usepackage{framed}
\makeatletter
\newenvironment{kframe}{%
 \def\at@end@of@kframe{}%
 \ifinner\ifhmode%
  \def\at@end@of@kframe{\end{minipage}}%
  \begin{minipage}{\columnwidth}%
 \fi\fi%
 \def\FrameCommand##1{\hskip\@totalleftmargin \hskip-\fboxsep
 \colorbox{shadecolor}{##1}\hskip-\fboxsep
     % There is no \\@totalrightmargin, so:
     \hskip-\linewidth \hskip-\@totalleftmargin \hskip\columnwidth}%
 \MakeFramed {\advance\hsize-\width
   \@totalleftmargin\z@ \linewidth\hsize
   \@setminipage}}%
 {\par\unskip\endMakeFramed%
 \at@end@of@kframe}
\makeatother

\definecolor{shadecolor}{rgb}{.97, .97, .97}
\definecolor{messagecolor}{rgb}{0, 0, 0}
\definecolor{warningcolor}{rgb}{1, 0, 1}
\definecolor{errorcolor}{rgb}{1, 0, 0}
\newenvironment{knitrout}{}{} % an empty environment to be redefined in TeX

\usepackage{alltt}

%%----------------------------------------------------------------------
%% opções comuns
\usepackage[brazilian]{babel}
\usepackage[utf8]{inputenc}
\usepackage[T1]{fontenc}
\usepackage{textcomp}
%\usepackage[margin=2cm]{geometry}
\usepackage{indentfirst}
\usepackage{fancybox}
%\usepackage[usenames,dvipsnames]{color}
\usepackage{amsmath,amsfonts,amssymb,amsthm}
\usepackage{lscape}
\usepackage{natbib}
\setlength{\bibsep}{0.0pt}
\usepackage{url}
\usepackage{multicol}
\usepackage{multirow}
\usepackage[final]{pdfpages}
\usepackage{setspace}
\usepackage{paralist} % enumitem, compactitem
  \plitemsep=2pt
%%----------------------------------------------------------------------

%%----------------------------------------------------------------------
%% FLOATS: graficos e tabelas
\usepackage{graphicx}
\usepackage{float} % fornece a opção [H] para floats
\usepackage{longtable}
\usepackage{supertabular}
%% captions e headings em sans-serif
\usepackage[font={sf},labelfont={sf,bf}]{caption}
\usepackage{subcaption}
\renewcommand{\thesubfigure}{\Alph{subfigure}}
\usepackage{titlesec}
\titleformat*{\section}{\normalsize\bfseries\sffamily}
\titleformat*{\subsection}{\normalsize\bfseries\sffamily}
\titleformat*{\subsubsection}{\normalsize\bfseries\sffamily}
\titleformat*{\paragraph}{\normalsize\bfseries\sffamily}
\titleformat*{\subparagraph}{\normalsize\bfseries\sffamily}
\theoremstyle{definition}
\newtheorem*{mydef}{Definição}
%%----------------------------------------------------------------------

%%----------------------------------------------------------------------
%% definiçoes de hyperref e xcolor
\usepackage{hyperref}
\usepackage{xcolor}
%%----------------------------------------------------------------------

%%----------------------------------------------------------------------
%% FONTES

%% micro-tipografia
\usepackage[protrusion=true,expansion=true]{microtype}
%% Bitstream Charter with mathdesign
\usepackage{lmodern} % sans-serif: Latin Modern
\usepackage[charter]{mathdesign} % serif: Bitstream Charter
\usepackage[scaled]{beramono} % truetype: Bistream Vera Sans Mono
\usepackage[scaled]{helvet}
%\usepackage{inconsolata}


%\usepackage[sf]{titlesec}
%%----------------------------------------------------------------------

%%----------------------------------------------------------------------
%% hifenização
\usepackage[htt]{hyphenat} % permite hifenizar texttt. Ao inves disso
% pode usar \allowbreak no ponto qu quiser quebrar dentro do texttt
\hyphenation{con-si-de-ra-ção pes-que-i-ros pes-que-i-ra se-gui-do-ras
  di-fe-ren-tes pla-ni-lha pla-ni-lhão re-fe-ren-te con-ta-gem
  em-bar-ques qua-li-da-de a-le-a-to-ri-za-dos}
%%----------------------------------------------------------------------

%%----------------------------------------------------------------------
%% comandos customizados
\usepackage{xspace} % lida com os espaços depois dos comandos
\providecommand{\eg}{\textit{e.g.}\xspace}
\providecommand{\ie}{\textit{i.e.}\xspace}
\providecommand{\R}{\textsf{R}\xspace}
\newcommand{\mb}[1]{\mathbf{#1}}
\newcommand{\bs}[1]{\boldsymbol{#1}}
\providecommand{\E}{\text{E}}
\providecommand{\Var}{\text{Var}}
\providecommand{\logit}{\text{logit}}
%% Para alterar o titulo do thebibliography
\addto\captionsbrazilian{%
  \renewcommand{\refname}{Bibliografia}
}
%%----------------------------------------------------------------------

%%----------------------------------------------------------------------
%% Comandos para deixar o texto mais compacto
\usepackage{marginnote}
\usepackage[top=1cm, bottom=1cm, inner=1cm, outer=1cm,nohead, nofoot, heightrounded, marginparsep=.05cm]{geometry}
\setlength{\parindent}{0pt}
%%----------------------------------------------------------------------
\IfFileExists{upquote.sty}{\usepackage{upquote}}{}
\begin{document}

\reversemarginpar % para colocar a marginnote a esquerda





\hrule
\vspace{0.3cm}

\begin{minipage}[c]{.85\textwidth}
  Bioestatística --- CE001 \\
  Prof. Fernando de Pol Mayer --- Departamento de Estatística --- DEST \\
  Exercícios: probabilidade \\
  Nome:   \hfill GRR: \hspace{2cm}
\end{minipage}\hfill
\begin{minipage}[c]{.15\textwidth}
\flushright
\includegraphics[width=2.2cm]{../img/ufpr-logo.png}
\end{minipage}

\vspace{0.3cm}
\hrule
\vspace{0.3cm}
%%----------------------------------------------------------------------

\begin{compactenum} % Magalhaes, pg 40
\item Para cada um dos eventos abaixo, escreva o espaço amostral
  correspondente e conte seus elementos:
  \begin{compactenum}
  \item Uma moeda é lançada duas vezes e observam-se as faces obtidas.
  \item Um dado é lançado duas vezes e a ocorrência de face par ou ímpar
    é observada.
  \item Uma urna contém 10 bolas azuis e 10 vermelhas. Três bolas são
    selecionadas ao acaso, com reposição, e as cores são anotadas.
  \item Dois dados são lançados simultaneamente e estamos interessados
    na soma das faces observadas.
  \item Em uma cidade famílias com 3 crianças são selecionadas ao acaso,
    anotando-se o sexo de cada uma, de acordo com a idade.
  \item Uma máquina produz 20 peças por hora, escolhe-se um instante
    qualquer e observa-se o número de defeituosas na próxima hora.
  \item Uma moeda é lançada consecutivamente até o aparecimento da
    primeira cara.
    % A partir daqui: Bussab, pg 108, cap 5
  \item Mede-se a duração de lâmpadas, deixando-as acesas até que se
    queimem.
  \item Um fichário com 10 nomes contém 3 nomes de mulheres.
    Seleciona-se ficha após ficha, até o último nome de mulher ser
    selecionado, e anota-se o número de fichas selecionadas.
  \item Uma moeda é lançada consecutivamente até o aparecimento da
    primeira cara e anota-se o número de lançamentos.
  \item De um grupo de 5 pessoas $\{A, B, C, D, E\}$, soteiam-se duas, uma
    após a outra, com reposição, e anota-se a configuração obtida.
  \item Mesmo enunciado anterior, mas sem reposição.
  \item Mesmo enunciado anterior, mas as duas selecionadas
    simultaneamente.
  \end{compactenum}

\vspace{0.3cm}
\hrule
\vspace{0.3cm}

% Bussab, pg 107, cap. 5
\item Uma urna contém duas bolas brancas e três bolas vermelhas.
  Retira-se uma bola ao acaso da urna. Se for branca lança-se uma moeda.
  Se for vermelha, ela é devolvida à urna e retira-se outra. Dê um
  espaço amostral para o experimento.

\vspace{0.3cm}
\hrule
\vspace{0.3cm}

% Ross, pg 71
\item Uma caixa contém 3 bolas de gude: 1 vermelha, 1 azul e 1
  branca. Considere um experimento que consiste em retirar duas bolas de
  gude desta caixa. Descreva o espaço amostral quando:
  \begin{compactenum}
  \item houver reposição da primeira bola retirada da caixa
  \item não houver reposição da primeira bola retirada da caixa
  \end{compactenum}

\vspace{0.3cm}
\hrule
\vspace{0.3cm}

% Montgomery, pg 19, cap 2
\item Uma balança digital é usada para fornecer pesos em gramas. Seja
  $A$ o evento em que um peso excede 11 gramas. Seja $B$ o evento em que
  um peso é menor que ou igual a 15 gramas, e seja $C$ o evento em que
  um peso é maior ou igual a 8 gramas e menor que 12 gramas. Descreva os
  seguintes eventos: \\
  \begin{inparaenum}
  \item $\Omega$ \,
  \item $A \cup B$ \,
  \item $A \cap B$ \,
  \item $A^{c}$ \,
  \item $A \cup B \cup C$ \,
  \item $(A \cup C)^{c}$ \,
  \item $A \cap B \cap C$ \,
  \item $B^{c} \cap C$ \,
  \item $A \cup (B \cap C)$
  \end{inparaenum}

\vspace{0.3cm}
\hrule
\vspace{0.3cm}

% Montgomery, pg 19, cap 2
\item Na fotossíntese, a qualidade da luz se refere aos comprimentos de
  onda de luz que são importantes. O comprimento de onda de uma amostra
  de radiações fotossinteticamente ativas (RFA) é medido em nanômetro. A
  faixa do vermelho é 675-700 nm, e a faixa do azul é 450-500 nm. Seja
  $A$ o evento em que RFA ocorre na faixa do vermelho, e $B$ o evento em
  que RFA ocorre na faixa do azul. Descreva o espaço amostral e indique
  cada um dos seguintes eventos: \\
  \begin{inparaenum}
  \item $A$ \,
  \item $B$ \,
  \item $A \cap B$ \,
  \item $A \cup B$
  \end{inparaenum}

\vspace{0.3cm}
\hrule
\vspace{0.3cm}

% Montgomery, pg 19, cap 2
\item Em replicação controlada, células são replicadas em um período de
  dois dias. DNA recém-sintetizado não pode ser replicado novamente até
  que a mtose seja completa. Dois mecanismos de controle foram
  identificados: um positivo e um negativo. Suponha que uma replicação
  seja observada em três células. Seja $A$ o evento em que todas as
  células são identificadas como positivas, e $B$ o evento em que todas
  as células são negativas. Descreva o espaço amostral e indique cada um
  dos seguintes eventos: \\
  \begin{inparaenum}
  \item $A$ \,
  \item $B$ \,
  \item $A \cap B$ \,
  \item $A \cup B$
  \end{inparaenum}

\vspace{0.3cm}
\hrule
\vspace{0.3cm}

\clearpage

\vspace{0.3cm}
\hrule
\vspace{0.3cm}

% Bussab, pg 113, cap. 5
\item Considere o lançamento de dois dados. Considere os eventos $A$ =
  ``soma dos números obtidos igual a 9'', e $B$ = ``número no primeiro
  dado maior ou igual a 4''.
  \begin{compactenum}
  \item Enumere os elementos de $A$ e $B$.
  \item Obtenha $A \cup B$, $A \cap B$, e $A^{c}$.
  \item Obtenha todas as probabilidades dos eventos acima.
  \end{compactenum}

\vspace{0.3cm}
\hrule
\vspace{0.3cm}

% Bussab, pg 126, cap. 5
\item Suponhamos que 10000 bilhetes sejam vendidos em uma loteria e 5000
  em outra, cada uma tendo apenas um ganhador. Um homem tem 100 bilhetes
  de cada. Qual a probabilidade de que:
  \begin{compactenum}
  \item Ele ganhe exatamente um prêmio?
  \item Ele ganhe alguma coisa?
  \end{compactenum}

\vspace{0.3cm}
\hrule
\vspace{0.3cm}

% Bussab, pg 126, cap. 5
\item Um grupo de 12 homens e 8 mulheres concorre a três prêmios através
  de um sorteio, sem reposição de seus nomes. Qual a probabilidade de:
  \begin{compactenum}
  \item Nenhum homem ser sorteado?
  \item Um prêmio ser ganho por homem?
  \item Dois homens serem premiados?
  \end{compactenum}

\vspace{0.3cm}
\hrule
\vspace{0.3cm}

% Ross pg 71
\item Suponha que $A$ e $B$ sejam eventos mutuamente exclusivos (ou
  disjuntos) para os quais $P(A) = 0,3$ e $P(B) = 0,5$. Qual é a
  probabilidade de que:
  \begin{compactenum}
  \item $A$ ou $B$ ocorra?
  \item $A$ ocorra mas $B$ não ocorra?
  \item $A$ e $B$ ocorram?
  \end{compactenum}

\vspace{0.3cm}
\hrule
\vspace{0.3cm}

% Montgomery, pg. 26, cap 2
\item Se $P(A) = 0,3$, $P(B) = 0,2$, e $P(A \cap B) = 0,1$. Determine as
  probabilidades: \\
    \begin{inparaenum}
    \item $P(A^c)$ \,
    \item $P(A \cup B)$ \,
    \item $P(A^c \cap B)$ \,
    \item $P(A \cap B^c)$ \,
    \item $P([A \cup B]^c)$ \,
    \item $P(A^c \cup B)$
    \end{inparaenum}

\vspace{0.3cm}
\hrule
\vspace{0.3cm}

% Montgomery, pg. 27, cap 2
\item Se $A$, $B$, e $C$ forem eventos mutuamente exclusivos, com $P(A)
  = 0,2$, $P(B) = 0,3$, e $P(C) = 0,4$. Determine as
  probabilidades: \\
    \begin{inparaenum}
    \item $P(A \cup B \cup C)$ \,
    \item $P(A \cap B \cap C)$ \,
    \item $P(A \cap B)$ \,
    \item $P([A \cup B] \cap C)$ \,
    \item $P(A^c \cap B^c \cap C^c)$
    \end{inparaenum}

\vspace{0.3cm}
\hrule
\vspace{0.3cm}

% Montgomery, pg 19, cap 2
\item Discos de plástico de policarbonato, provenientes de um fornecedor,
  são analisados com relação à resistência a arranhões e a choques. Os
  resultados de uma amostra de 100 discos estão resumidos a seguir:
  \begin{table}[!h]
    \centering
    \begin{tabular}{lcc}
      \hline
      \multirow{2}{*}{\textbf{Res. a arranhões}}
      & \multicolumn{2}{l}{\textbf{Res. a choques}} \\
      \cline{2-3}
                & Alta         & Baixa       \\
      \hline
      Alta      & 70           & 9           \\
      Baixa     & 16           & 5           \\
      \hline
    \end{tabular}
  \end{table}

  Seja $A$ o evento em que um disco tem alta resistência a choque e $B$
  o evento em que um disco tem alta resistência a arranhões. Com isso:
  \begin{compactenum}
  \item Determine o número de discos em $A \cap B$, $A^{c}$, e $A \cup B$.
  \item Se um disco for selecionado aleatoriamente, determine as
    seguintes probabilidades: \\
    \begin{inparaenum}
    \item $P(A)$ \,
    \item $P(B)$ \,
    \item $P(A^{c})$ \,
    \item $P(A \cap B)$ \,
    \item $P(A \cup B)$ \,
    \item $P(A^{c} \cup B)$ \,
    \item $P(A|B)$ \,
    \item $P(B|A)$
    \end{inparaenum}
  \item Se um disco for selecionado ao acaso, qual será a probabilidade
    de sua resistência a arranhões ser alta e de sua resistência a
    choque ser alta?
  \item Se um disco for selecionado ao acaso, qual será a probabilidade
    de sua resistência a arranhões ser alta ou de sua resistência a
    choque ser alta?
  \item Os eventos $A$ e $B$ são mutuamente exclusivos?
  \item Os eventos $A$ e $B$ são independentes?
  \end{compactenum}

\vspace{0.3cm}
\hrule
\vspace{0.3cm}

% magalhaes pg 41
\item Sejam $A$ e $B$ dois eventos em um espaço amostral, tais que $P(A)
  = 0,2$, $P(B) = p$, $P(A \cup B) = 0,5$, e $P(A \cap B) =
  0,1$. Determine o valor de $p$.

\vspace{0.3cm}
\hrule
\vspace{0.3cm}

\clearpage

\vspace{0.3cm}
\hrule
\vspace{0.3cm}

% Montgomery, pg 31, cap 2
\item Um lote de 100 chips semicondutores contém 20 defeituosos. Dois
  deles são selecionados ao acaso, sem reposição.
  \begin{compactenum}
  \item Qual é a probabilidade de que o primeiro chip selecionado seja
    defeituoso?
  \item Qual é a probabilidade de que o segundo chip selecionado seja
    defeituoso, dado que o primeiro deles foi defeituoso?
  \item Qual é a probabilidade de que ambos sejam defeituosos?
  \item Como a resposta do item (b) mudaria se os chips selecionados
    fossem repostos antes da próxima seleção?
  \end{compactenum}

\vspace{0.3cm}
\hrule
\vspace{0.3cm}

% Montgomery, pg 21, cap 2
\item A tabela abaixo resume 204 reações endotérmicas envolvendo
  bicarbonato de sódio.
  \begin{table}[!h]
    \centering
    \begin{tabular}{ccc}
      \hline
      \multirow{2}{*}{\textbf{Condições finais de temperatura}}
      & \multicolumn{2}{c}{\textbf{Calor absorvido}} \\
      \cline{2-3}
                & Abaixo do valor alvo         & Acima do valor alvo \\
      \hline
      266 K      & 12           & 40           \\
      271 K      & 44           & 16           \\
      274 K      & 56           & 36           \\
      \hline
    \end{tabular}
  \end{table}

  Seja $A$ o evento em que a temperatura final de uma reação seja 271 K
  ou menos. Seja $B$ o evento em que o calor absorvido esteja acima do
  valor alvo. Com isso:
  \begin{compactenum}
  \item Determine o número de reações em cada um dos seguintes eventos: \\
    \begin{inparaenum}
    \item $A \cap B$ \,
    \item $A^c$ \,
    \item $A \cup B$ \,
    \item $A \cup B^c$ \,
    \item $A^c \cap B^C$
    \end{inparaenum}
  \item Determine as  seguintes probabilidades: \\
    \begin{inparaenum}
    \item $P(A \cap B)$ \,
    \item $P(A^c)$ \,
    \item $P(A \cup B)$ \,
    \item $P(A \cup B^c)$ \,
    \item $P(A^c \cap B^c)$ \,
    \item $P(A^c \cup B^c)$ \,
    \item $P(A|B)$ \,
    \item $P(A^c|B)$ \,
    \item $P(A|B^c)$ \,
    \item $P(B|A)$
    \end{inparaenum}
  \item Os eventos $A$ e $B$ são independentes?
  \end{compactenum}

\vspace{0.3cm}
\hrule
\vspace{0.3cm}

% Montgomery pg 33, cap 2
\item Suponha que $P(A|B) = 0,4$ e $P(B) = 0,5$. Determine o seguinte: \\
  \begin{inparaenum}
  \item $P(A \cap B)$ \,
  \item $P(A^c \cap B)$
  \end{inparaenum}

\vspace{0.3cm}
\hrule
\vspace{0.3cm}

% Montgomery pg 33, cap 2
\item Suponha que $P(A|B) = 0,2$, $P(A|B^c) = 0,3$ e $P(B) = 0,8$. Qual
  é $P(A)$? (Dica: escreva $A$ como a união de dois eventos disjuntos).

\vspace{0.3cm}
\hrule
\vspace{0.3cm}


% Montgomery, pg 21, cap 2
\item Um artigo na revista \textit{The Journal of Data Science},
  forneceu a seguinte tabela de falhas em poços, para grupos de
  diferentes formações geológicas em Baltimore (EUA):
  \begin{table}[!h]
    \centering
    \begin{tabular}{lcc}
      \hline
      \multirow{2}{*}{\textbf{Grupo com formação geológica}}
      & \multicolumn{2}{c}{\textbf{Poços}} \\
      \cline{2-3}
                & Falha      & Total  \\
      \hline
      Gnaise     & 170           & 1685           \\
      Granito    & 2           & 28           \\
      Mina Loch de xisto    & 443           & 3733   \\
      Máfico & 14 & 363 \\
      Mármore & 29 & 309 \\
      Mina Prettyboy de xisto & 60 & 1403 \\
      Outros xistos & 46 & 933 \\
      Serpentina & 3 & 39 \\
      \hline
    \end{tabular}
  \end{table}

  Seja $A$ o evento em que a formação geológica tenha mais de 1000 poços
  e $B$ o evento em que o poço tenha falhado. Com isso:
  \begin{compactenum}
  \item Determine o número de poços dos seguintes eventos: \\
    \begin{inparaenum}
    \item $A \cap B$ \,
    \item $A^c$ \,
    \item $A \cup B$ \,
    \item $A \cup B^c$ \,
    \item $A^c \cap B^C$
    \end{inparaenum}
  \item Determine as seguintes probabilidades: \\
    \begin{inparaenum}
    \item $P(A \cap B)$ \,
    \item $P(A^c)$ \,
    \item $P(A \cup B)$ \,
    \item $P(A \cup B^c)$ \,
    \item $P(A^c \cap B^c)$ \,
    \item $P(A^c \cup B^c)$ \,
    \item $P(A|B)$
    \end{inparaenum}
  \item Qual a probabilidade de uma falha, dado que existem mais de 1000
    falhas em uma formação geológica?
  \item Qual a probabilidade de uma falha, dado que existem menos de 500
    falhas em uma formação geológica?
  \item Os eventos $A$ e $B$ são independentes?
  \end{compactenum}

\vspace{0.3cm}
\hrule
\vspace{0.3cm}

\clearpage

\vspace{0.3cm}
\hrule
\vspace{0.3cm}

% Montgomery, pg 20, cap 2
\item O tempo de enchimento de um reator é medido em minutos (e frações
  de minutos). Seja $\Omega = \mathbb{R}^{+}$. Defina os aventos $A$ e
  $B$ como segue:
  \begin{equation*}
    A = \{x : x \leq 72,5\} \qquad \text{e} \qquad B = \{x : x > 52,5\}
  \end{equation*}
  Descreva cada um dos seguintes eventos: \\
  \begin{inparaenum}
  \item $A^{c}$ \,
  \item $B^{c}$ \,
  \item $A \cap B$ \,
  \item $A \cup B$
  \end{inparaenum}

\vspace{0.3cm}
\hrule
\vspace{0.3cm}

% Montgomery, pg 33, cap 2
\item Falhas no coração são por causa tanto de ocorrências naturais
  (87\%) como por fatores externos (13\%). Fatores externos estão
  relacionados a substâncias induzidas (73\%) ou a objetos estranhos
  (27\%). Ocorrências naturais são causadas por bloqueio arterial
  (56\%), doenças (27\%) e infecção (17\%).
  \begin{compactenum}
  \item Determine a probabilidade de uma falha ser causada por
    substância induzida.
  \item Determine a probabilidade de uma falha ser causada por doença ou
    infecção.
  \end{compactenum}

\vspace{0.3cm}
\hrule
\vspace{0.3cm}

% Montgomery, pg 20, cap 2
\item Uma amostra de dois itens é selecionada sem reposição a partir de
  uma batelada. Descreva o espaço amostral (ordenado) para cada uma das
  seguintes bateladas:
  \begin{compactenum}
  \item A batelada contém os itens $\{a, b, c, d\}$
  \item A batelada contém os itens $\{a, b, c, d, e, f, g\}$
  \item A batelada contém 4 itens defeituosos e 20 itens bons
  \item A batelada contém 1 item defeituoso e 20 itens bons
  \end{compactenum}

\vspace{0.3cm}
\hrule
\vspace{0.3cm}

% Montgomery, pg 22, cap 2
\item Cada um dos cinco resultados possíveis de um experimento aleatório
  é igualmente provável. O espeço amostral é $\Omega = \{a, b, c, d,
  e\}$. Seja $A$ o evento $\{a, b\}$ e $B$ o evento $\{c, d, e\}$.
  Determine: \\
  \begin{inparaenum}
  \item $P(A)$ \,
  \item $P(B)$ \,
  \item $P(A^c)$ \,
  \item $P(A \cup B)$ \,
  \item $P(A \cap B)$
  \end{inparaenum}

\vspace{0.3cm}
\hrule
\vspace{0.3cm}

% Montgomery, pg 22, cap 2
\item O espaço amostral de um experimento aleatório é $\Omega = \{a, b, c, d,
  e\}$, com probabilidades 0,1; 0,1; 0,2; 0,4; 0,2, respectivamente.
  Seja $A$ o evento $\{a, b, c\}$ e $B$ o evento $\{c, d, e\}$.
  Determine: \\
  \begin{inparaenum}
  \item $P(A)$ \,
  \item $P(B)$ \,
  \item $P(A^c)$ \,
  \item $P(A \cup B)$ \,
  \item $P(A \cap B)$
  \end{inparaenum}

\vspace{0.3cm}
\hrule
\vspace{0.3cm}


% Montgomery, pg 20, cap 2
\item Uma amostra de duas placas de circuito impresso é selecionada sem
  reposição a partir de uma batelada. Descreva o espaço amostral
  (ordenado) para cada uma das seguintes bateladas:
  \begin{compactenum}
  \item A batelada contém 90 placas que são não defeituosas, 8 placas
    com pequenos defeitos, e 2 placas com grandes defeitos.
  \item A batelada contém 90 placas que são não defeituosas, 8 placas
    com pequenos defeitos, e 1 placa com grandes defeitos.
  \end{compactenum}

\vspace{0.3cm}
\hrule
\vspace{0.3cm}

% Montgomery, pg. 24, cap. 2
\item Em uma titulação ácido-base, uma base ou um ácido é gradualmente
  adicionada(o) ao outro até que eles sejam completamente neutralizados.
  Uma vez que ácidos e bases são geralmente incolores, o pH é medido
  para monitorar a reação. Suponha que o ponto de equivalência seja
  alcançado depois que aproximadamente 100 ml de uma solução de NaOH
  tenham sido adicionados (o suficiente para reagir com todo o ácido
  acético presente), porém essa quantidade pode variar de 95 ml a 104
  ml. Suponha que volumes sejam medidos em ml em uma escala discreta, e
  descreva o espaço amostral.
  \begin{compactenum}
  \item Qual é a probabilidade de que a equivalência seja indicada em
    100 ml?
  \item Qual é a probabilidade de que a equivalência seja indicada em
    menos do que 100 ml?
  \item Qual é a probabilidade de que a equivalência seja indicada entre
    98 ml e 102 ml (inclusive)?
  \item Considere que dois técnicos conduzam a titulação de forma
    independente.
    \begin{compactenum}
    \item Qual é a probabilidade de ambos os técnicos obterem
      equivalência em 100 ml?
    \item Qual é a probabilidade de ambos os técnicos obterem
      equivalência entre 98 e 104 ml (inclusive)?
    \end{compactenum}
  \end{compactenum}

\vspace{0.3cm}
\hrule
\vspace{0.3cm}

% Ross pg 71
\item Uma loja aceita cartões de crédito American Express ou Visa. Um
  total de 24\% de seus consumidores possui um cartão American Express,
  61\% possuem Visa, e 11\% possuem ambos. Que percentual desses
  consumidores possui um cartão aceito pelo estabelecimento?

\vspace{0.3cm}
\hrule
\vspace{0.3cm}

\clearpage

\vspace{0.3cm}
\hrule
\vspace{0.3cm}

% Montgomery, pg. 24, cap. 2
\item Em uma bateria de NiCd,  uma célula completamente cerregada é
  composta de Hidróxido de Níquel. Níquel é um elemento que tem
  múltiplos estados de oxidação, sendo geralmente encontrado nos
  seguintes estados:
  \begin{table}[!h]
    \centering
    \begin{tabular}{cc}
      \hline
      \textbf{Carga de níquel}
      & \textbf{Proporções encontradas} \\
      \hline
      0 & 0,17 \\
      +2 & 0,35 \\
      +3 & 0,33 \\
      +4 & 0,15 \\
      \hline
    \end{tabular}
  \end{table}
  \begin{compactenum}
  \item Qual é a probabilidade de uma célula ter no mínimo uma das
    opções de níquel carregado positivamente?
  \item Qual é a probabilidade de uma célula não ser composta de uma
    carga positiva de níquel maior do que +3?
  \end{compactenum}

\vspace{0.3cm}
\hrule
\vspace{0.3cm}

% Magalhaes, pg 41
\item Uma Universidade tem 10000 alunos, dos quais 4000 são considerados
  esportistas. Temos ainda que 500 alunos são do curso de Economia, 700
  são de Administração, 100 são esportistas e da Economia, e 200 são
  esportistas e da Administração. (Os restantes dos alunos que não se
  encaixam em nenhum curso podem ser colocados na categoria
  ``outros''). Monte a tabela de contingência, e calcule as
  probabilidades de um aluno selecionado ao acaso:
  \begin{compactenum}
  \item Ser esportista
  \item Ser esportista e aluno da Administração
  \item Não ser da Economia nem da Administração
  \item Não ser esportista, e ser de outros cursos
  \item Ser esportista, dado que faz Economia
  \item Fazer Administração, dado que não é esportista
  \item Não ser esportista, dado que não faz nem Economia, nem
    Administração
  \item Fazer outros cursos, dado que é esportista
  \end{compactenum}

\vspace{0.3cm}
\hrule
\vspace{0.3cm}

% Kazmier, pg 78
\item Suponhamos que um um aluno estime que sua probabilidade de receber um
  conceito final ``A'' em Estatística é 0,6, e a probabilidade de um
  ``B'' é de 0,4 (ele não considera receber menos do que isso!). Com isso:
  \begin{compactenum}
  \item Os eventos ``receber A'' e ``receber B'' são mutuamente
    exclusivos? Por que?
  \item Determine a probabilidade condicional de que obtenha um ``B'',
    dado que de fato tenha recebido um ``A''.
  \item Verifique se os eventos ``receber A'' e ``receber B'' são
    independentes.
  \end{compactenum}

\vspace{0.3cm}
\hrule
\vspace{0.3cm}

% morettin, pg 23
\item Sejam $A$ e $B$ eventos tais que $P(A) = 0,2$, $P(B) = p$, $P(A
  \cup B) = 0,6$. Calcular o valor de $p$ considerando $A$ e $B$
  \begin{compactenum}
  \item mutuamente exclusivos
  \item independentes
  \end{compactenum}

\vspace{0.3cm}
\hrule
\vspace{0.3cm}

% morettin, pg 40
\item Uma urna contém 10 bolas verdes e 6 azuis. Tiram-se duas bolas ao
  acaso, sem reposição. Qual a probabilidade de que as duas bolas:
  \begin{compactenum}
  \item sejam verdes?
  \item sejam da mesma cor?
  \item sejam de cores diferentes?
  \end{compactenum}

\vspace{0.3cm}
\hrule
\vspace{0.3cm}

% kazmier, pg 77
\item De 100 pessoas que solicitaram emprego de programador de
  computadores durante um ano em uma grande empresa, 40
  possuíam experiência anterior ($E$), e 30 possuíam certificado
  profissional ($C$). Vinte dos candidatos possuíam tanto experiência
  anterior como certificado profissional e foram incluídos nas contagens
  dos dois grupos.
  \begin{compactenum}
  \item Elaborar um diagrama de Venn para representar estes eventos.
  \item Qual a probabilidade de que um candidato escolhido ao acaso
    tenha experiência ou certificado?
  \item Qual a probabilidade de que um candidato escolhido ao acaso
    tenha experiência ou certificado, mas não ambos?
  \item Determinar a probabilidade condicional de que um candidato
    escolhido ao acaso tenha um certificado, dado que ele tenha alguma
    experiência anterior.
  \item Verificar se os eventos $E$ e $C$ são independentes.
  \end{compactenum}

\vspace{0.3cm}
\hrule
\vspace{0.3cm}

\clearpage

\vspace{0.3cm}
\hrule
\vspace{0.3cm}

% kazmier, pg 78
\item Em geral, a probabilidade de que um possível cliente faça uma
  compra quando procurado por um vendedor é de 0,4. Se um vendedor
  seleciona do arquivo, aleatoriamente, três clientes, e faz contato com
  os mesmos, qual a probabilidade de que os três façam compras?

\vspace{0.3cm}
\hrule
\vspace{0.3cm}

% Montgomery, pg 31, cap 2
\item Um artigo na revista \textit{The Canadian Entomologist} estudou a
  vida da praga da alfafa a partir dos ovos até a vida adulta. A tabela
  seguinte mostra o número de larvas que sobreviveram em cada estágio do
  desenvolvimento.
  \begin{table}[!h]
    \centering
    \begin{tabular}{lp{2cm}p{2cm}lll}
      Ovos & Fase precoce da larva & Fase madura da larva & Pré-pupa
      & Pupa & Adultos \\
      \hline
      421 & 412 & 306 & 45 & 35 & 31 \\
    \end{tabular}
  \end{table}
  \begin{compactenum}
  \item Qual é a probabilidade de um ovo sobreviver até a vida adulta?
  \item Qual é a probabilidade de sobrevivência até a vida adulta, dada
    a sobrevivência para a fase madura da larva?
  \item Que estágio tem a menor probabilidade de sobrevivência para o
    próximo estágio?
  \end{compactenum}

\vspace{0.3cm}
\hrule
\vspace{0.3cm}

% Montgomery, pg 36, cap 2
\item Se $P(A|B) = 0,4$, $P(B) = 0,8$, $P(A) = 0,5$, os eventos $A$ e
  $B$ são independentes?

\vspace{0.3cm}
\hrule
\vspace{0.3cm}

% Montgomery, pg 36, cap 2
\item Se $P(A|B) = 0,3$, $P(B) = 0,8$, $P(A) = 0,3$, o evento $B$ e o
  evento complementar de $A$ são independentes?

\vspace{0.3cm}
\hrule
\vspace{0.3cm}

% Montgomery, pg 36, cap 2
\item Se $P(A) = 0,2$, $P(B) = 0,2$, e $A$ e $B$ são mutuamente
  excludentes, eles são independentes?

\vspace{0.3cm}
\hrule
\vspace{0.3cm}

% Montgomery, pg 36, cap 2
\item Matriz redundante de discos independentes (RAID -
  \textit{Redundant Array of Independent Disks}) é uma tecnologia que
  usa discos rígidos múltiplos para aumentar a velocidade de
  transferência de dados e fornecer cópia de segurança instantânea de
  dados. Suponha que a probabilidade de qualquer disco rígido falhar em
  um dia seja 0,001, e que as falhas do disco sejam independentes.
  \begin{compactenum}
  \item Suponha que você implemente um esquema de RAID 0, que usa dois
    discos rígidos, cada um contendo uma imagem do outro, como um
    espelho. Qual é a probabilidade de perda de dados? Considere que a
    perda de dados ocorrerá se ambos os discos falharem dentro do mesmo
    dia.
  \item Suponha que você implemente um esquema de RAID 1, que divide os
    dados em dois discos rígidos. Qual é a probabilidade de perda de
    dados? Considere que a perda de dados ocorrerá se no mínimo um disco
    falhar dentro do mesmo dia. (Dica: escreva o evento ``no mínimo um
    disco falhar'' como o seu complementar).
  \end{compactenum}

\vspace{0.3cm}
\hrule
\vspace{0.3cm}

% Montgomery, pg 44, cap 2
\item Cabelos vermelhos naturais consistem em dois genes. Pessoas com
  cabelo vermelho natural têm dois genes dominantes, dois genes
  recessivos, ou um dominante e um recessivo. Um grupo de 1000 pessoas
  foi categorizado como segue:
  \begin{table}[!h]
    \centering
    \begin{tabular}{lccc}
      \hline
      \multirow{2}{*}{\textbf{Gene 1}}
      & \multicolumn{3}{c}{\textbf{Gene 2}} \\
      \cline{2-4}
                & Dominante      & Recessivo & Outro  \\
      \hline
      Dominante     & 5           & 25 & 30           \\
      Recessivo    & 7           & 63 & 35           \\
      Outro    & 20           & 15 & 800   \\
      \hline
    \end{tabular}
  \end{table}

  Seja $A$ o evento em que uma pessoa tem um gene dominante de cabelo
  vermelho, e seja $B$ o evento em que uma pessoa tem um gene recessivo
  de cabelo vermelho. Se uma pessoa desse grupo for selecionada ao acso,
  calcule o seguinte:
  \begin{compactenum}
  \item $P(A)$
  \item $P(A \cap B)$
  \item $P(A \cup B)$
  \item $P(A^c \cap B)$
  \item $P(A|B)$
  \item Considerando que para uma pessoa ter cabelo vermelho são
    necessários dois genes dominates, qual a probabilidade de que a
    pessoa selecionada tenha cabelo vermelho?
  \end{compactenum}

\vspace{0.3cm}
\hrule
\vspace{0.3cm}


\end{compactenum}

\end{document}
